\def\topic{Continuous integral kernels}
\ifdefined\MyHeaderIncluded
\else
\def\qwertyflag{true}

\ifdefined\MyHeaderIncluded

\else
  \newcommand{\foobar}{FooBar}
  \def\MyHeaderIncluded{true}
% if is till the end of the file

\documentclass[10pt,a4paper]{article}
\usepackage[utf8]{inputenc}
%\usepackage[german]{babel}
%\usepackage[cyrillic]{babel}
\usepackage[T1]{fontenc}
%\usepackage[T2A]{fontenc}
%\usepackage[english, russian]{babel}
\usepackage{amsmath}
\usepackage{amsfonts}
\usepackage{amssymb}
\usepackage{amsthm}
\usepackage{mathtools} %coloneqq

\usepackage{bussproofs} %

%\renewcommand{\thesubsection}{\arabic{subsection}}
%\newtheorem{thm}{Theorem}[section]
%\renewcommand{\thethm}{\arabic{subsection}.\arabic{thm}}

\usepackage{thmtools}
\usepackage{hyperref}
\hypersetup{
    colorlinks=true,
    linkcolor=blue,
    filecolor=magenta,      
    urlcolor=cyan,
    pdftitle={Overleaf Example},
    %pdfpagemode=FullScreen,
    }
%\usepackage{stix-type1}%\turnediota
\theoremstyle{definition}

\newtheorem{constructor}{Constructor}[section]
\newtheorem{iaxiom}{Axiom}
\newtheorem{axiom}{Axiom}
\newtheorem*{axiom*}{Axiom}
\newtheorem{definition}{Definition}[section]
\newtheorem*{definition*}{Definition}
\newtheorem{remark}{Remark}[section]
\newtheorem*{remark*}{Remark}
\newtheorem{recall}{Recall}[section]
\newtheorem*{recall*}{Recall}

\newtheorem{lemma}{Lemma}[section]
\newtheorem{theorem}{Theorem}[section]
\newtheorem*{theorem*}{Theorem}
\newtheorem{note}{Note}[section]
\usepackage{etoolbox}

%\newcommand*{\@Dashv}[2]{%
%  \reflectbox{$\m@th#1#2$}%
%}
\newcommand*{\separation}[3]{% 1)binder, 2)set, 3)condition
  \{#1\in #2: #3\}
}

% https://www.freecodecamp.org/news/simple-regex-tricks-for-beginners-3acb3fa257cb/
%  \\mathcal\{T\} |---> \topology
%  \\mathcal\{([^\}]*)\}   |--->  $1
\newcommand{\AND}{\ \ \&\ \ }
\newcommand{\any}{\mathrm{\ any}}
\newcommand{\borel}{\mathcal{B}}
\newcommand{\can}{\mathrm{can}}
\newcommand{\canonize}{\mathrm{canonize}}
\newcommand{\closure}{\mbox{cl}}
\newcommand{\complex}{\mathbb{C}}
\newcommand{\converges}{\longrightarrow}
\newcommand{\countable}{\mbox{\ countable}}
\newcommand{\ldefined}{\mbox{\ defined}}
\newcommand{\dom}{\mbox{dom}}
\newcommand{\Dom}{\mbox{Dom}}
\newcommand{\Done}{\noindent \textbf{Done.}\\}
\newcommand{\lexists}{\mbox{\ exists}}
\newcommand{\lforall}{\mbox{\ forall}}
\newcommand{\lunique}{\mbox{\ unique}}
\newcommand{\family}{\mathrm{\ family}}
\newcommand{\Fm}{\mathrm{Fm}}
\newcommand{\FuncSymb}{\mathrm{FuncSymb}}
\newcommand{\function}{\mathrm{\ function}}
\newcommand{\functional}{\mbox{\ functional}}
\newcommand{\IFF}{\ \Longleftrightarrow\ }
\renewcommand{\iff}{\ \longleftrightarrow\ }
\newcommand{\IMPLIES}{\ \Longrightarrow\ }
\renewcommand{\implies}{\ \longrightarrow\ }
\newcommand{\Inf}{\mathrm{Inf}}
\newcommand{\injective}{\mbox{\ injective}}
\newcommand{\injection}{\mbox{\ injection}}
\newcommand{\interior}{\mbox{int}}
\newcommand{\intersection}{\bigcap}
\newcommand{\isempty}{\mathrm{\ empty}}
\newcommand{\LB}{\mathrm{LB}}
\newcommand{\nghbset}{\mathcal{N}}
\newcommand{\nonempty}{\mathrm{\ nonempty}}
\newcommand{\Ord}{\mathrm{Ord}}
\newcommand{\power}{\mathcal{P}}
\newcommand{\pr}{\mathrm{pr}}
\newcommand{\fst}{\mathrm{fst}}
\newcommand{\snd}{\mathrm{snd}}
\newcommand{\PredSymb}{\mathrm{PredSymb}}
\newcommand{\Proof}{\noindent \textbf{Proof.}\\}
\newcommand{\ran}{\mathrm{ran}}
\newcommand{\Ran}{\mathrm{Ran}}
\newcommand{\real}{\mathbb{R}}
\newcommand{\relation}{\mathrm{\ relation}}
\newcommand{\set}{\mathrm{\ set}}
\newcommand{\singleton}{\mathrm{\ singleton}}
\newcommand{\semantics}{\mathrm{semantics}}
\newcommand{\semialgebra}{\mathcal{S}}
\newcommand{\Sup}{\mathrm{Sup}}
\newcommand{\surjective}{\mbox{\ surjective\ on\ }}
\newcommand{\Term}{\mathrm{Term}}
\newcommand{\Type}{\mathrm{Type}}
\newcommand{\topology}{\mathcal{T}}
\newcommand{\total}{\mbox{\ total\ on\ }}
\newcommand{\transitive}{\mathrm{\ transitive}}
\newcommand{\union}{\bigcup}
\newcommand{\lTrue}{\mathrm{True}}
\newcommand{\lFalse}{\mathrm{False}}
\newcommand{\lUncertain}{\mathrm{Uncertain}}
\newcommand{\lif}{\mathrm{if}}
\newcommand{\lthen}{\mathrm{then}}
\newcommand{\lelse}{\mathrm{else}}
\newcommand{\lexception}{\mathrm{exception}}
\newcommand{\lundef}{\mathrm{undef}}
\newcommand{\univocal}{\mathrm{\ univocal}}
\newcommand{\univocalat}{\mathrm{\ univocal\ at\ }}
\newcommand{\UB}{\mathrm{UB}}
\newcommand{\Un}{\mathcal{U}}
\newcommand{\Var}{\mathrm{Var}}

\newcommand{\notvarphi}{\phi}
\renewcommand{\phi}{\varphi}

\newcommand{\comment}[1]{\quad\mbox{(#1)}}
% local commands
%\preto\section{\undef\blipo} %etoolbox
\preto\subsection{\undef\blipo} %etoolbox
% https://tex.stackexchange.com/questions/105684/limit-macro-newcommand-scope-to-section
\usepackage{graphicx}
\usepackage{authblk}

\setlength\parindent{0pt} % whole doc of \noindent
\usepackage{soul}
%\usepackage{mathgl}

\fi % guard
\ifdefined\topic
\else
\def\topic{No Titel}
\fi
\title{\topic}
\author[1]{Georgy Dunaev}
\begin{document}
\maketitle
\fi
\section{A criterion for separately continuous kernels}
\label{sec:continuous_kernel_criterion}
The criterion below proves a locally finite Radon-measure version of
Simon's Theorem~A.1.3 \parencite{Simon1982Semigroups}, with boundedness of
the kernel explicit. The source uses Baire measures, asks for a
norm-continuous and weakly continuous extension, and imposes full
support for uniqueness. We retain the existence criterion without
full support, and state precisely where support is needed. Its
conclusion is separate continuity, not joint continuity. We give the
proof, including the dual-pair argument needed on a noncompact space.

\subsection{Measures, representatives, and the dual pair}
Let $(X,d)$ be a nonempty separable locally compact metric space, and let
$\nu$ be a positive locally finite Radon measure on $X$. Write
$C_b(X)$ for the complex bounded continuous functions, with supremum
norm, and $\mathcal M(X)$ for the finite complex Radon measures, with
total variation norm. A function $f\in L^1(X,\nu)$ determines the
measure $f\nu\in\mathcal M(X)$, with
$\|f\nu\|_{\mathrm{TV}}=\|f\|_1$. Local compactness and separability
make $X$ a countable union of compact sets; in particular $\nu$ is
$\sigma$-finite. The next lemma justifies the Radon property of the
finite density measures, including for the wider measure convention
discussed below. All measures here are on the Borel sets. In this
metric setting the Baire and Borel sigma algebras coincide: each closed
set is the zero set of its continuous distance function.

\begin{lemma}[Regularity of the finite measures used here]
\label{lem:kernel_finite_measure_regularity}
Every finite positive Borel measure $\mu$ on $X$ is outer regular
by open sets and inner regular by compact sets. Consequently every
finite complex Borel measure on $X$ is a finite Radon measure.
\end{lemma}
\begin{proof}
Separability and metrizability give a countable base. Local compactness
then gives a countable cover by relatively compact open sets: choose
a relatively compact neighborhood at each point and extract a
countable subcover using the base. Finite unions of their closures
form increasing compact sets $K_n$ with union $X$.

We first establish approximation by closed and open sets. Let
$\mathcal R$ be the class of Borel sets $E$ for which, for every
$\epsilon>0$, there are a closed $F$ and an open $G$ with
$F\subseteq E\subseteq G$ and $\mu(G\setminus F)<\epsilon$.
Every open set $G\ne X$ belongs to $\mathcal R$: the closed sets
$\{x:d(x,X\setminus G)\geq1/n\}$ increase to $G$, and continuity
from below gives the required approximation. The set $X$ itself is
immediate. The class $\mathcal R$ is closed under complements by
interchanging the complements of $F,G$.

For $E=\bigcup_{j\geq1}E_j$ with $E_j\in\mathcal R$, choose
$F_j\subseteq E_j\subseteq G_j$ with
$\mu(G_j\setminus F_j)<\epsilon 2^{-j}/3$. Since $\mu$ is finite,
choose $N$ with
$\mu(E\setminus\bigcup_{j\leq N}E_j)<\epsilon/3$.
Then $F=\bigcup_{j\leq N}F_j$ is closed,
$G=\bigcup_{j\geq1}G_j$ is open, and
\[
\begin{aligned}
\mu(G\setminus F)
&\leq\sum_{j\geq1}\mu(G_j\setminus E_j)
 +\mu\left(E\setminus\bigcup_{j\leq N}E_j\right)\\
&\quad +\sum_{j\leq N}\mu(E_j\setminus F_j)<\epsilon.
\end{aligned}
\]
Thus $\mathcal R$ is a sigma algebra containing the open sets, so it
contains every Borel set. For a closed inner approximant $F$,
$F\cap K_n$ is compact and increases to $F$; continuity from below
therefore improves closed inner approximation to compact inner
approximation. The same argument applied to the finite positive
total variation of a complex measure proves the final assertion.
\end{proof}

Use the bilinear pairing
\[
\langle\mu,g\rangle=\int_X g\,d\mu,
\qquad \mu\in\mathcal M(X),\quad g\in C_b(X).
\]
For $x\in X$, the Dirac measure $\delta_x$ is the unit point mass at
$x$, so $\langle\delta_x,g\rangle=g(x)$.
Here $\sigma(\mathcal M,C_b)$ means convergence against every fixed
$g\in C_b$, and $\sigma(C_b,\mathcal M)$ means convergence against
every fixed $\mu\in\mathcal M$. These definitions apply to nets,
not only to sequences. They specify a dual pair; they do not identify
$\mathcal M(X)$ with the full Banach dual of $C_b(X)$, an identification
that is generally false on a noncompact space.
For completeness, the pairing separates both spaces. Dirac measures
separate continuous functions. If a finite complex measure integrates
every bounded continuous function to zero, apply this to
$\max(0,1-n d(x,F))$ for each nonempty closed $F$. Dominated convergence
gives $\mu(F)=0$. Closed inner approximation with respect to $|\mu|$
then gives $\mu(E)=0$ for every Borel $E$, so $\mu=0$.

Let $J:C_b(X)\to L^\infty(X,\nu)$ take the almost-everywhere class.
A bounded operator $T:L^1(\nu)\to L^\infty(\nu)$ has an extension
$S:\mathcal M(X)\to C_b(X)$ in the sense used here if
\[
J(S(f\nu))=Tf\qquad(f\in L^1(\nu)).
\]
This is an equality of equivalence classes. Full support means that
every nonempty open subset of $X$ has positive $\nu$-measure.
If $\nu$ has full support,
$J$ is injective and isometric: a continuous function exceeding a bound
at one point exceeds it on an open set of positive measure. Without
full support, continuous representatives need not be unique.

\begin{lemma}[Continuous linear functionals for the specified weak topology]
\label{lem:kernel_weak_dual}
Every continuous complex-linear functional $F$ on
$(\mathcal M(X),\sigma(\mathcal M,C_b))$ has the form
$F(\mu)=\int g\,d\mu$ for a unique $g\in C_b(X)$.
\end{lemma}
\begin{proof}
Continuity at zero gives $g_1,\ldots,g_n\in C_b(X)$ and
$\epsilon>0$ such that $|\int g_j\,d\mu|<\epsilon$ for all $j$
implies $|F(\mu)|<1$. If all those integrals vanish, every scalar
multiple of $\mu$ belongs to this neighborhood; linearity forces
$F(\mu)=0$. Thus $F$ factors through the map
$\mu\mapsto(\int g_1\,d\mu,\ldots,\int g_n\,d\mu)$ into
$\mathbb C^n$. A linear functional on its image extends to
$\mathbb C^n$ by extending a basis, so it is a linear combination
of its coordinates. The same combination of the $g_j$ gives $g$.
Testing Dirac measures $\delta_x$ proves uniqueness.
\end{proof}

\subsection{The criterion and its proof}
\begin{theorem}[Bounded separately continuous kernel criterion]
\label{thm:separate_kernel_criterion}
Let $T:L^1(X,\nu)\to L^\infty(X,\nu)$ be bounded and linear.
The following are equivalent:
\begin{enumerate}
\item There is a bounded function $K:X\times X\to\mathbb C$,
continuous in each variable separately, such that
\[
(Tf)(x)=\int_X K(x,y)f(y)\,d\nu(y)
\]
for $\nu$-almost every $x$, for every $f\in L^1(\nu)$.
\item There is a bounded linear extension
$S:\mathcal M(X)\to C_b(X)$ of $T$ which is continuous from
$\sigma(\mathcal M,C_b)$ to $\sigma(C_b,\mathcal M)$.
\end{enumerate}
In particular these conditions imply $\operatorname{Ran}T\subseteq
J(C_b(X))$. Kernels and extensions correspond by
\[
K(x,y)=(S\delta_y)(x),\qquad
(S\mu)(x)=\int_X K(x,y)\,d\mu(y),
\]
and $\|S\|=\sup_{x,y}|K(x,y)|$.
If $\nu$ has full support, both the kernel and the extension are
unique, and their norm is also $\|T\|_{L^1\to L^\infty}$.
\end{theorem}
\begin{proof}
First suppose $K$ satisfies condition~1, and put
$B=\sup_{x,y}|K(x,y)|$. The function $K$ is jointly Borel measurable,
although it need not be jointly continuous. To verify measurability,
choose a countable dense sequence $(a_j)$ in $X$. For each positive
integer $n$, let $q_n(x)$ be the first $a_j$ with
$d(x,a_j)<1/n$. Its level sets are Borel. Consequently
$K(q_n(x),y)$ is Borel on $X\times X$, since it is given on a
countable Borel partition in $x$ by continuous functions of $y$.
Continuity in the first variable gives its pointwise limit $K(x,y)$.

Define $S$ by the asserted integral formula. It is linear and
$|(S\mu)(x)|\le B\|\mu\|_{\mathrm{TV}}$. If $x_n\to x$,
separate continuity and dominated convergence with respect to $|\mu|$
give $(S\mu)(x_n)\to(S\mu)(x)$. Since $X$ is metric, this proves
$S\mu\in C_b(X)$. The kernel representation of $T$ proves the
extension property.

For $\lambda\in\mathcal M(X)$, define
\[
g_\lambda(y)=\int_X K(x,y)\,d\lambda(x).
\]
The same dominated-convergence argument shows $g_\lambda\in C_b(X)$.
Fubini applies to finite total variation measures, because $K$ is
bounded and jointly measurable, and gives
\[
\langle\lambda,S\mu\rangle
=\int_X g_\lambda\,d\mu.
\]
Each defining scalar test of the target weak topology is therefore
a continuous scalar test of the source weak topology. This proves
the required continuity for nets, not merely sequential continuity.

Conversely, suppose condition~2 holds and define
$K(x,y)=(S\delta_y)(x)$. For fixed $y$, this is continuous in $x$
because $S\delta_y\in C_b(X)$. For fixed $x$, the map
$F_x:\mu\mapsto(S\mu)(x)$ is continuous for
$\sigma(\mathcal M,C_b)$: evaluation at $x$ is the target test
$\delta_x$, and $S$ is weakly continuous. By
Lemma~\ref{lem:kernel_weak_dual}, there is $g_x\in C_b(X)$ such that
$F_x(\mu)=\int g_x\,d\mu$. Evaluating at $\delta_y$ identifies
$g_x(y)=K(x,y)$. This proves continuity in $y$ and, for every measure,
the asserted integral formula for $S$. Moreover
$|K(x,y)|\le\|S\|$, since $\|\delta_y\|_{\mathrm{TV}}=1$.
Putting $\mu=f\nu$ and using the extension property proves condition~1.
The formula bounds $\|S\|$ by $\sup|K|$, while Dirac measures give
the reverse inequality. The two constructions are inverse, since
the formula for $S\mu$ was derived for all measures.

Now assume full support. For fixed $y\in X$, choose relatively
compact open neighborhoods $U_n$ of $y$, contained in the radius
$1/n$ ball. Then $0<\nu(U_n)<\infty$. Put
$f_n=\mathbf1_{U_n}/\nu(U_n)$, so $\|f_n\|_1=1$.
The continuous representative $S(f_n\nu)$ has supremum norm
$\|Tf_n\|_\infty\le\|T\|$. For fixed $x$, continuity of $K(x,\cdot)$
gives
\[
(S(f_n\nu))(x)=\frac1{\nu(U_n)}\int_{U_n}K(x,u)\,d\nu(u)
\longrightarrow K(x,y).
\]
Hence $|K(x,y)|\le\|T\|$, and integration gives the reverse
operator-norm inequality, proving norm equality.

If $K_1,K_2$ represent the same $T$, the functions
$\int(K_1(x,y)-K_2(x,y))f(y)\,d\nu(y)$ are continuous in $x$
and vanish almost everywhere, so full support makes them zero at
every $x$. Substitute the same $f_n$ and take the limit to get
$K_1(x,y)=K_2(x,y)$ for every $x,y$. The already proved integral
formula then gives equality of the corresponding extensions.
\end{proof}

\paragraph{The measure convention and a wider formulation.}
The source allows a Baire measure which need not have finite total
mass. Our locally finite Radon convention already includes Lebesgue
measure on $\mathbb R^d$. In fact the bounded-kernel equivalence above
holds for every positive Borel measure $\nu$ on the same space $X$,
without local finiteness: each $f\in L^1(\nu)$ defines a finite complex
Borel measure $f\nu$, and finite Borel measures on a separable locally
compact metric space are Radon. The proof integrates against these
finite measures, so its dominated-convergence and Fubini arguments
remain valid; it does not need a product of two possibly non-sigma-finite
copies of $\nu$.

For the uniqueness and norm-equality conclusions it suffices that
every nonempty open set $U$ contain a Borel set $E$ with
$0<\nu(E)<\infty$. Replace $U_n$ in the proof by such a subset
$E_n\subset U_n$ and use $\mathbf1_{E_n}/\nu(E_n)$; continuity
still makes its averages tend to the value at the center. This
condition holds for locally finite measures of full support, and more
generally for full-support semifinite measures. Here semifiniteness
means that each measurable set of infinite measure contains a measurable
subset of finite positive measure.

A measure convention is necessary for an unrestricted full-support
statement. The Borel measure assigning $\nu(\varnothing)=0$ and
$\nu(E)=\infty$ to every nonempty $E$ has full support but
$L^1(\nu)=\{0\}$. Both $K=0$ and $K=1$ then represent its zero
operator. They yield different bounded weakly continuous extensions,
namely $S\mu=0$ and $S\mu=\mu(X)\mathbf1$. Thus full support alone
does not ensure uniqueness for arbitrary nonsemifinite Borel measures.
This explains the measure hypotheses rather than assigning a broader
meaning to the source's Baire-measure convention.

\subsection{Why weak continuity is part of the hypothesis}
Norm boundedness of an extension alone is insufficient, even when the
reference measure has full support. On $X=\mathbb R$ with Lebesgue
measure, let $T:L^1\to L^\infty$ be zero and define
\[
(S\mu)(x)=\mu(\{0\})\qquad(x\in\mathbb R).
\]
This is a bounded linear map $\mathcal M(\mathbb R)\to C_b(\mathbb R)$
of norm one. It extends $T$, since every density $f\,dx$ assigns zero
mass to the singleton. Nevertheless it differs from the zero extension.

To see exactly which hypothesis fails, put
\[
d\mu_n(x)=\frac n2\mathbf1_{(-1/n,1/n)}(x)\,dx.
\]
For every $g\in C_b(\mathbb R)$,
\[
\left|\int g\,d\mu_n-g(0)\right|
\leq\sup_{|x|<1/n}|g(x)-g(0)|\longrightarrow0.
\]
Thus $\mu_n\to\delta_0$ in $\sigma(\mathcal M,C_b)$, whereas
$S\mu_n=0$ and $S\delta_0=\mathbf1$. Evaluation at zero is a test
in the target weak topology, so $S$ fails the required weak continuity.

This approximation is not convergence in total variation:
$\|\mu_n-\delta_0\|_{\mathrm{TV}}=2$. The partition into $\{0\}$
and its complement gives the lower bound two, and the triangle
inequality gives the upper bound. Consequently the uniqueness of the
extension in the theorem comes from the specified weak topology,
rather than from a nonexistent norm density of $L^1$ densities among
all finite measures.

\subsection{What the criterion does and does not assert}
No full-support assumption was used for the equivalence. It was used
only for uniqueness and comparison with the norm of $T$.
For example, let $X=\{0,1\}$ with the discrete topology and
$\nu=\delta_0$. Both $K=0$ and the kernel which is one at $(1,1)$
and zero elsewhere represent the zero operator on $L^1(\nu)$,
but their extensions to $\mathcal M(X)$ differ. In particular the
extension norm cannot be recovered from $\|T\|$ in this setting.
Boundedness of the kernel in the theorem is explicit: with
$X=\mathbb R$, $\nu=\delta_0$, the unbounded separately continuous
function $K(x,y)=y$ also represents the zero operator, but it does
not define a bounded extension on all finite measures. Statements
about arbitrary representatives outside the support require care.

Separate continuity cannot be strengthened to joint continuity by
this argument. On $X=\mathbb R$ with Lebesgue measure, set
\[
K(x,y)=
\begin{cases}
xy/(x^2+y^2),&(x,y)\ne(0,0),\\
0,&(x,y)=(0,0).
\end{cases}
\]
This kernel is separately continuous and bounded by $1/2$, but
$K(t,t)=1/2$ for $t\ne0$, so it is not jointly continuous at the
origin. It defines a bounded $L^1\to L^\infty$ operator and a
weakly continuous extension exactly as in the theorem. Full support
makes this its unique separately continuous kernel; there is no
alternative jointly continuous kernel for the same operator.
Obtaining joint continuity for a heat semigroup therefore requires
an additional argument using that semigroup's analytic structure.

\ifdefined\qwertyflag
\undef\qwertyflag
\end{document}
\fi
